%=============================================================================!
% 18.8  TRAINING, VALIDATION AND TEST                                         !
%=============================================================================!
\section{Training, validation and test}
\label{sec:lo-splits}

A student who revises from past papers and then sits an exam made of the
same questions scores well without having learned the subject; the
score measures memory alone. A fitted model is judged the
same way, on readings it has not seen, and the judgement is only as
honest as the separation of those readings from the ones it learned
from. With the frames of a molecular dynamics run that separation is
subtler than it looks.

\subsection*{Three sets}

The readings are divided before any fitting. The \emph{training} set
fixes the weights; the \emph{validation} set chooses everything that the
training set cannot, such as the degree of a polynomial, the penalty
$\lambda$, the learning rate or the moment to stop
(\cref{sec:lo-ridge}); the \emph{test} set is opened once, at the end, to
report how well the chosen model predicts. Every choice made by looking
at a set spends some of that set's honesty, which is why the test is
touched last.

The division assumes that the sets are independent samples of the
readings the model will meet. Frames of a run are not: a frame
\SI{10}{\femto\second} after another holds nearly the same positions and
nearly the same energy, and \cref{sec:cv-correlation} measured how long
such a memory lasts, as the statistical inefficiency $g$, the number of
consecutive samples worth one independent one.

\subsection*{Frames that leak}

Two runs of Chapter~12's liquid argon were made with every
\SI{10}{\femto\second} step kept, as first-principles runs often are, and
a model was asked to learn each frame's potential energy $U$ from the
histogram of its \num{32640} pair distances in 110 bins of
\SI{0.05}{\angstrom}. The task is well posed: $U$ is the sum of the pair
energy over the pairs, so it is linear in the histogram, through the pair
energy at each bin, with a miss from the binning alone of
\SI{22.5}{\milli\electronvolt} against a spread of $U$ of
\SI{175}{\milli\electronvolt}. Ridge regression ($\lambda = 10^{-2}$, the
110 counts scaled by the training frames' mean and spread) learned $U$
from the first frames of one run, with a fifth held out for validation,
and was tested on frames of the other run, every \SI{50}{\femto\second},
a run it never saw (\cref{fig:lo-leak}).

At \SI{10}{\femto\second} the energy has $g = 109$ frames, so a training
run of \SI{2}{\pico\second}, 200 frames, holds about two independent
ones. Held out at random, its validation frames each have training
frames \SI{10}{\femto\second} either side of them, nearly identical, and
the validation miss is $(30.5 \pm 4.4)$\,\si{\milli\electronvolt} over ten
random splits, while the independent run gives $(128 \pm
16)$\,\si{\milli\electronvolt}, four times more: the random split measures
how well the model recalls its neighbours. Held out in whole blocks of
\SI{0.5}{\pico\second}, the validation miss is $(92 \pm
22)$\,\si{\milli\electronvolt}, close to the truth. With \SI{4}{\pico
\second} the random split still reports 21.7 against 37.8, the blocked
one 34.8; only at 10 and \SI{20}{\pico\second}, about 9 and 18 independent
frames, do all three agree near the binning's own floor, since a model
with enough independent frames no longer needs the neighbours.

The cure is to split by time, into whole blocks longer than the memory
$g$, or better whole runs, for validation and test. The same care
applies to every quantity computed from the training set alone, such as
the mean and spread that scale the features, which must not be computed
with the validation frames in. A learned potential's reported error
means little without the split that produced it.

\begin{figure}[htb]
\centering
\includegraphics{ch18_learning/leak}
\caption{Consecutive frames leak between splits. Liquid argon, a frame
every \SI{10}{\femto\second}; ridge regression learns each frame's
potential energy from the histogram of its pair distances. (a) The
energy over \SI{4}{\pico\second}, with the frames that a random split
holds out (dots) and those that a split into blocks of
\SI{0.5}{\pico\second} holds out (shaded). (b) The rms miss of $U$ on
validation frames held out at random and in blocks, and on an
independent run, against the length of the training run, with the
spread over ten splits.}
\label{fig:lo-leak}
\end{figure}

\begin{keyidea}
Training readings fix the weights, validation readings choose the
settings, and test readings, touched once, report the result. All three
must be independent samples; consecutive frames of a run are not, and
split at random they let a model look up its neighbours, reporting a miss
four times too small here. Split by blocks longer than the statistical
inefficiency, or by whole runs.
\end{keyidea}

\notebook{18}{split a run's frames at random and in blocks and compare the
misses}

\begin{selfcheck}
A first-principles run of \SI{10}{\pico\second} keeps every
\SI{1}{\femto\second} step, and its energy has $g = 400$ steps. About how
many independent frames does it hold, and how long should each validation
block be?
\end{selfcheck}
